% !TeX root = ../DECO_buku_iot_cerdas.tex
\chapter{KESIMPULAN DAN SARAN}

\section{Kesimpulan}

Berdasarkan hasil perancangan dan pengujian sistem monitoring suhu berbasis
IoT cerdas, diperoleh beberapa kesimpulan sebagai berikut:

\begin{enumerate}
    \item Sistem monitoring suhu berbasis IoT menggunakan ESP32 dan
    DS18B20 berhasil dirancang dan disimulasikan menggunakan Wokwi. Sensor
    DS18B20 berhasil membaca suhu lingkungan dan mengirimkan data ke
    ThingSpeak melalui koneksi Wi-Fi virtual Wokwi setiap 15 detik.
    \item ESP32 berhasil membaca data suhu dari sensor DS18B20 dan
    mengirimkannya ke ThingSpeak melalui koneksi Wi-Fi virtual Wokwi
    menggunakan protokol HTTP API.
    \item ThingSpeak berhasil digunakan sebagai platform dashboard untuk
    menampilkan data suhu aktual (Field~1), prediksi suhu 1 menit ke depan
    (Field~2), dan status AI (Field~3) secara otomatis dan berkala.
    \item Model \textit{Random Forest Regressor} berhasil diintegrasikan
    menggunakan Python untuk membaca data historis ThingSpeak dan memprediksi
    suhu 1 menit ke depan dengan tujuh fitur \textit{time series} berbasis lag,
    perubahan suhu, dan rata-rata bergerak.
    \item Sistem yang dikembangkan telah memenuhi konsep IoT cerdas karena
    mencakup akuisisi data sensor, komunikasi \textit{cloud}, dashboard
    monitoring, serta integrasi AI untuk analisis prediktif secara
    \textit{end-to-end}.
\end{enumerate}

\section{Saran}

Beberapa saran pengembangan sistem ke depan adalah:

\begin{enumerate}
    \item Menambahkan sensor tambahan seperti DHT22, MQ2, atau LDR agar data
    masukan AI lebih bervariasi dan model dapat mempelajari korelasi antar
    parameter lingkungan.
    \item Menggunakan data historis yang lebih panjang agar performa prediksi
    \textit{Random Forest} lebih stabil dan representatif terhadap pola suhu
    yang sebenarnya.
    \item Menambahkan indikator teks atau notifikasi pada dashboard agar
    status AI tidak hanya berupa kode numerik, sehingga lebih informatif
    bagi pengguna.
    \item Mengembangkan sistem pada perangkat ESP32 fisik agar tidak hanya
    berjalan pada simulasi Wokwi, sehingga pengukuran suhu bersifat nyata
    dan responsif terhadap perubahan lingkungan.
    \item Menambahkan aktuator seperti \textit{buzzer}, LED, atau \textit{relay}
    apabila sistem ingin dikembangkan menjadi sistem peringatan otomatis
    berbasis prediksi suhu.
    \item Mengeksplorasi model AI lain seperti LSTM (\textit{Long Short-Term
    Memory}) untuk membandingkan performa prediksi \textit{time series}
    dengan \textit{Random Forest Regressor}.
\end{enumerate}
